\par\noindent\textbf{(1)} 沿特征 \(dx/dt=u\)，\(du/dt=0\)，故初始标记 \(\xi\) 上 \(u(x,t)=u_0(\xi)\)、\(x=\xi+t u_0(\xi)\)。在映射 \(\xi\mapsto x\) 单调时该式唯一；若出现特征相交，单值经典解失效。
\par\noindent\textbf{(2)} 对隐式映射求导 \(\partial\xi/\partial x=[1+t u_0'(\xi)]^{-1}\)，所以 \(u_x=u_0'(\xi)/[1+t u_0'(\xi)]\)。若 \(\min_\xi u_0'(\xi)<0\)，梯度首次发散时刻 \(t_*=-1/\min u_0'\)，是有限时间波形破碎。
\par\noindent\textbf{(3)} 方程可写 \(u_t+[u^2/2]_x=0\)。对左右状态 \(u_L,u_R\) 的跳跃面 \(x=s t\) 积分守恒律得 \(s(u_L-u_R)=u_L^2/2-u_R^2/2\)，故 \(s=(u_L+u_R)/2\)。这是守恒跃变条件，不保证任意跳跃都物理可接受。
\par\noindent\textbf{(4)} 通量 \(f(u)=u^2/2\) 为凸函数，粘性极限给Lax熵条件 \(u_L>s>u_R\) 对应 \(u_L>u_R\) 的压缩激波。若 \(u_L<u_R\)，物理解为扇形稀疏波 \(u=x/t\) 在两态之间，不得用不满足熵条件的膨胀跳跃替代。
